% --- Άσκηση 35033 (35033.tex) ---
\Askhsh{35033}{2}
\begin{ylh}{2}
    \item 2.2. Διάταξη Πραγματικών Αριθμών
    \item 2.3. Απόλυτη Τιμή Πραγματικού Αριθμού
    \item 3.1. Εξισώσεις 1ου Βαθμού
\end{ylh}
Δίνονται οι παραστάσεις $A = |2x - 4|\ \text{και}\ B = |x - 3|$.

\begin{alist}
\item Αν $2 < x < 3,$ να δείξετε ότι $A = 2x - 4$.
\monades{7}
\item Αν $2 < x < 3$, να δείξετε ότι $A + B = x - 1$.
\monades{9}
\item Υπάρχει $x \in (2,\left. \ 3 \right)$ ώστε να ισχύει $A + B = 2;$ Να αιτιολογήσετε την απάντησή σας.
\monades{9}
\end{alist}
